% ========================================================
% BÀI TẬP TOÁN 5 - KẾT NỐI TRI THỨC VỚI CUỘC SỐNG
% Bài 32: Ôn tập một số hình phẳng
% Định dạng: BTTeX Part III - Trắc nghiệm điền số ngắn
% ========================================================

% Câu 1
\begin{ex}
  Hình tam giác có bao nhiêu đỉnh?
  \shortans{3}
  \loigiai{
    Tam giác có 3 đỉnh.
  }
\end{ex}

% Câu 2
\begin{ex}
  Hình tứ giác có bao nhiêu cạnh?
  \shortans{4}
  \loigiai{
    Tứ giác có 4 cạnh.
  }
\end{ex}

% Câu 3
\begin{ex}
  Hình thang có mấy cạnh đáy song song với nhau?
  \shortans{2}
  \loigiai{
    Hình thang có 2 cạnh đáy song song.
  }
\end{ex}

% Câu 4
\begin{ex}
  Hình thoi có mấy cặp cạnh đối diện song song?
  \shortans{2}
  \loigiai{
    Hình thoi có 2 cặp cạnh đối diện song song.
  }
\end{ex}

% Câu 5
\begin{ex}
  Trong một hình chữ nhật có bao nhiêu góc vuông?
  \shortans{4}
  \loigiai{
    Hình chữ nhật có 4 góc vuông.
  }
\end{ex}

% Câu 6
\begin{ex}
  Hình vuông có bao nhiêu trục đối xứng?
  \shortans{4}
  \loigiai{
    Hình vuông có 4 trục đối xứng.
  }
\end{ex}

% Câu 7
\begin{ex}
  Hình chữ nhật có bao nhiêu trục đối xứng?
  \shortans{2}
  \loigiai{
    Hình chữ nhật có 2 trục đối xứng.
  }
\end{ex}

% Câu 8
\begin{ex}
  Một tam giác đều có chu vi $45\text{ cm}$. Độ dài mỗi cạnh là bao nhiêu xăng-ti-mét?
  \shortans{15}
  \loigiai{
    $45 : 3 = 15\text{ cm}$.
  }
\end{ex}

% Câu 9
\begin{ex}
  Một hình vuông có chu vi $60\text{ cm}$. Độ dài cạnh của hình vuông là bao nhiêu xăng-ti-mét?
  \shortans{15}
  \loigiai{
    $60 : 4 = 15\text{ cm}$.
  }
\end{ex}

% Câu 10
\begin{ex}
  Một hình thoi có chu vi $80\text{ cm}$. Độ dài mỗi cạnh của hình thoi là bao nhiêu xăng-ti-mét?
  \shortans{20}
  \loigiai{
    $80 : 4 = 20\text{ cm}$.
  }
\end{ex}

% Câu 11
\begin{ex}
  Một hình tam giác cân có cạnh bên dài $12\text{ cm}$, cạnh đáy dài $10\text{ cm}$. Chu vi hình tam giác cân đó là bao nhiêu xăng-ti-mét?
  \shortans{34}
  \loigiai{
    $12 \times 2 + 10 = 34\text{ cm}$.
  }
\end{ex}

% Câu 12
\begin{ex}
  Một hình bình hành có độ dài hai cạnh liên tiếp là $12\text{ cm}$ và $8\text{ cm}$. Chu vi hình bình hành đó là bao nhiêu xăng-ti-mét?
  \shortans{40}
  \loigiai{
    $(12 + 8) \times 2 = 40\text{ cm}$.
  }
\end{ex}

% Câu 13
\begin{ex}
  Đường kính của hình tròn có bán kính $15\text{ cm}$ là bao nhiêu xăng-ti-mét?
  \shortans{30}
  \loigiai{
    $15 \times 2 = 30\text{ cm}$.
  }
\end{ex}

% Câu 14
\begin{ex}
  Bán kính của hình tròn có đường kính $36\text{ cm}$ là bao nhiêu xăng-ti-mét?
  \shortans{18}
  \loigiai{
    $36 : 2 = 18\text{ cm}$.
  }
\end{ex}

% Câu 15
\begin{ex}
  Tam giác vuông có một góc vuông $90^\circ$ và một góc nhọn $35^\circ$. Góc nhọn còn lại có số đo là bao nhiêu độ?
  \shortans{55}
  \loigiai{
    $90^\circ - 35^\circ = 55^\circ$.
  }
\end{ex}

% Câu 16
\begin{ex}
  Tổng số đo ba góc trong một tam giác bằng bao nhiêu độ?
  \shortans{180}
  \loigiai{
    Tổng 3 góc bằng $180^\circ$.
  }
\end{ex}

% Câu 17
\begin{ex}
  Một hình thang cân có hai góc kề đáy lớn mỗi góc bằng $70^\circ$. Tổng số đo hai góc kề đáy bé là bao nhiêu độ?
  \shortans{220}
  \loigiai{
    Mỗi góc đáy bé là $180 - 70 = 110^\circ$. Tổng là $110 \times 2 = 220^\circ$.
  }
\end{ex}

% Câu 18
\begin{ex}
  Một hình thoi có độ dài hai đường chéo là $12\text{ cm}$ và $16\text{ cm}$. Hai đường chéo cắt nhau tại trung điểm của mỗi đường và tạo thành bao nhiêu góc vuông?
  \shortans{4}
  \loigiai{
    Hai đường chéo vuông góc tạo thành 4 góc vuông.
  }
\end{ex}

% Câu 19
\begin{ex}
  Một hình lục giác đều được chia thành bao nhiêu hình tam giác đều nhỏ bằng nhau bằng các đường chéo chính?
  \shortans{6}
  \loigiai{
    Được chia thành 6 hình tam giác đều.
  }
\end{ex}

% Câu 20
\begin{ex}
  Một hình chữ nhật có nửa chu vi là $25\text{ cm}$, chiều dài hơn chiều rộng $5\text{ cm}$. Chiều dài là bao nhiêu xăng-ti-mét?
  \shortans{15}
  \loigiai{
    $(25 + 5) : 2 = 15\text{ cm}$.
  }
\end{ex}

% Câu 21
\begin{ex}
  Ở câu trên, chiều rộng là bao nhiêu xăng-ti-mét?
  \shortans{10}
  \loigiai{
    $25 - 15 = 10\text{ cm}$.
  }
\end{ex}

% Câu 22
\begin{ex}
  Một hình tam giác có chiều cao là $8\text{ cm}$, độ dài đáy gấp đôi chiều cao. Đáy tam giác dài bao nhiêu xăng-ti-mét?
  \shortans{16}
  \loigiai{
    $8 \times 2 = 16\text{ cm}$.
  }
\end{ex}

% Câu 23
\begin{ex}
  Một tứ giác có độ dài 4 cạnh lần lượt là $5\text{ cm}, 7\text{ cm}, 8\text{ cm}, 10\text{ cm}$. Chu vi tứ giác đó là bao nhiêu xăng-ti-mét?
  \shortans{30}
  \loigiai{
    $5 + 7 + 8 + 10 = 30\text{ cm}$.
  }
\end{ex}

% Câu 24
\begin{ex}
  Cần bao nhiêu điểm phân biệt không thẳng hàng để tạo thành một tam giác?
  \shortans{3}
  \loigiai{
    Cần 3 điểm.
  }
\end{ex}

% Câu 25
\begin{ex}
  Một hình thang có độ dài hai đáy là $14\text{ cm}$ và $6\text{ cm}$. Trung bình cộng hai đáy là bao nhiêu xăng-ti-mét?
  \shortans{10}
  \loigiai{
    $(14 + 6) : 2 = 10\text{ cm}$.
  }
\end{ex}
